\documentclass[11pt]{article}
\usepackage[utf8]{inputenc}
\usepackage{amsmath,amssymb}
\usepackage[margin=1in]{geometry}
\usepackage{hyperref}
\emergencystretch=3em

\title{The $d=2$ theorem for the chart amplitude $\eta e^{\beta s\xi}$}
\author{Claude, with Daniel Murfet}
\date{2026-09-07}

\begin{document}
\maketitle
\sloppy

\begin{abstract}
The amplitude-level hypotheses of unit~88 (axis Lipschitz bounds and the mixed second difference)
are derived for $\Phi(u,v,s)=\eta(u,v)e^{\beta s\xi(u,v)}$ from Lipschitz and mixed-difference bounds
on $\xi,\eta$ themselves, which hold for $C^2$ functions. The rectangle difference of a product,
$D[fg]=D[f]g_3+(f_1-f_0)(g_3-g_1)+(f_2-f_0)(g_3-g_2)+f_0D[g]$, and the identity
$e^{t_3}-e^{t_1}-e^{t_2}+1=(e^{t_1}-1)(e^{t_2}-1)+(e^{t_3}-e^{t_1+t_2})$ give
$|D[\eta e^{\beta s\xi}]|\le C_2uv(1+s)^2e^{5\beta Ls}$. Consequently the two-dimensional chart integral
with equal exponents satisfies $Z(N)=N^{-p}(A\log N+B)+O(N^{-(p+\delta)}(1+\log N))$ with
$A=(k_1k_2)^{-1}\eta(0,0)S_p(\xi(0,0))$; the $n$-form follows with $\tfrac A2\log n$. Zero
\texttt{sorry}/\texttt{axiom}.
\end{abstract}

\section{The things formalised}
{\small
\begin{verbatim}
theorem amp_sub_le (β L M ξ₁ ξ₂ η₁ η₂ s : ℝ) (hβ : 0 < β) (hs : 0 ≤ s) (hξ₁ : |ξ₁| ≤ L)
    (hξ₂ : |ξ₂| ≤ L) (hη₂ : |η₂| ≤ M) :
    |η₁ * Real.exp (β * s * ξ₁) - η₂ * Real.exp (β * s * ξ₂)|
      ≤ (|η₁ - η₂| + M * (β * s) * |ξ₁ - ξ₂|) * Real.exp (β * s * L)
theorem amp_mixed_le ... :
    |e₃ * Real.exp (β * s * a₃) - e₁ * Real.exp (β * s * a₁) - e₂ * Real.exp (β * s * a₂)
        + e₀ * Real.exp (β * s * a₀)|
      ≤ (K + 2 * K ^ 2 * β + M * β ^ 2 * K ^ 2 + M * β * K) * (u * v) * ((1 + s) ^ 2
        * Real.exp (5 * β * L * s))
noncomputable def chartAmp (β : ℝ) (ξ η : ℝ → ℝ → ℝ) (u v s : ℝ) : ℝ :=
  η u v * Real.exp (β * s * ξ u v)
theorem twoDChart_second_order ... (hξu hξv hηu hηv : Lipschitz in each variable)
    (hξD hηD : mixed second differences ≤ K * (u * v)) :
    (fun N : ℝ => twoDAmp β b N h₁ h₂ k₁ k₂ (chartAmp β ξ η)
        - N ^ (-p) * (twoDA β p k₁ k₂ (chartAmp β ξ η) * Real.log N
          + twoDB β b p k₁ k₂ (chartAmp β ξ η))) =O[atTop]
      fun N : ℝ => N ^ (-(p + min (k₁ : ℝ)⁻¹ (k₂ : ℝ)⁻¹)) * (1 + Real.log N)
theorem twoDChart_A (β p : ℝ) (k₁ k₂ : ℕ) (ξ η : ℝ → ℝ → ℝ) :
    twoDA β p k₁ k₂ (chartAmp β ξ η)
      = 1 / ((k₁ : ℝ) * k₂) * (η 0 0 * weightedMass β (ξ 0 0) (p - 1))
\end{verbatim}
}

\section{Roadblocks and resolutions}
\begin{itemize}
\item \texttt{linarith} needs the product $\beta sL\ge0$ as a named hypothesis to compare
$2\beta sL$ with $4\beta Ls$ (nonlinear monomials are atoms).
\item A single-step \texttt{calc} as the last tactic of a \texttt{have} block swallowed the following
\texttt{exact} line as a further step; use \texttt{exact} directly for one-step bounds.
\end{itemize}

\section{Outcome}
The $d=2$ log polynomial is available for the paper's chart integral with $\xi,\eta$ of class $C^2$
(via Lipschitz and mixed-difference bounds). Next: the face-parameter second-order theorem under the
strict gap, the $C\ne0$ variants of the leading theorem, and a third Astra consult on all-order
subtraction.
\end{document}
